\newpage
\phantomsection
\label{prf:infimum-less-than-supremum}

\begin{remark*}[Return]
\hyperref[thm:infimum-less-than-supremum]{Return to Theorem}
\end{remark*}

\begin{theorem*}[Infimum-Supremum Comparison]
For any non-empty bounded set $A \subseteq \mathbb{R}$, the infimum of $A$ is less than or equal to the supremum of $A$.
\end{theorem*}

\begin{proof}
\textbf{Professional Standard Proof.}~\\
Let $A$ be a non-empty bounded set in $\mathbb{R}$. By the completeness of $\mathbb{R}$, $i = \inf A$ and $s = \sup A$ exist. Since $A \neq \varnothing$, choose $a \in A$. By the definition of infimum, $i \le a$. By the definition of supremum, $a \le s$. By the transitivity of $\le$, $i \le a \le s$, which implies $i \le s$.
\end{proof}

\begin{proof}
\textbf{Detailed Learning Proof.}~\\

\textbf{Step 1. Invoke Completeness.}
Because $A$ is bounded and non-empty, the Least Upper Bound Property (and its counterpart for infima) ensures that $s = \sup A$ and $i = \inf A$ are well-defined real numbers.

\textbf{Step 2. Utilize Non-emptiness.}
The assumption $A \neq \varnothing$ is necessary. It allows us to assert the existence of at least one element $a \in A$.

\textbf{Step 3. Apply the Lower Bound property.}
By definition, the infimum is a lower bound for $A$. Therefore, for the specific element $a$ chosen in Step 2:
\[ i \le a \]

\textbf{Step 4. Apply the Upper Bound property.}
By definition, the supremum is an upper bound for $A$. Therefore, for the same element $a$:
\[ a \le s \]

\textbf{Step 5. Chain by Transitivity.}
Combining the results of Step 3 and Step 4:
\[ i \le a \le s \]
By the transitivity of the standard order on $\mathbb{R}$, we conclude $i \le s$.
\end{proof}

\begin{remark*}[Proof structure]
This is a direct proof utilizing a single element of the set as a logical pivot between the lower bound property of the infimum and the upper bound property of the supremum.
\end{remark*}

\begin{remark*}[Dependencies]~\\
\begin{itemize}
  \item \hyperref[pred:nonempty-set]{Definition of NonemptySet}
  \item \hyperref[pred:infimum]{Definition of Infimum (as a LowerBound)}
  \item \hyperref[pred:supremum]{Definition of Supremum (as an UpperBound)}
  \item \hyperref[pred:least-upper-bound-property]{Least Upper Bound Property of $\mathbb{R}$}
\end{itemize}
\end{remark*}